\section{Structured-mask disintegration of the physical columns}
\label{sec:tpc33-structured-mask}

We now prove the physical reduction announced in
\eqref{eq:tpc33-intro-affine-column}.  The result uses the actual mask
\eqref{eq:tpc33-actual-mask}.  It is false as a statement about every
entrywise bounded mask.

\subsection{An exact projection of the row-gcd mask}

For \(G\ge1\), define
\begin{equation}
 \lambda_G(v)
 :=\sum_{\substack{g\mid v\\g\le G}}\mu(v/g).
 \label{eq:tpc33-truncated-gcd-coefficient}
\end{equation}

\begin{lemma}[Exact row-gcd projector]
\label{lem:tpc33-row-gcd-projector}
For every \(d,e\in\N\),
\begin{equation}
 \boxed{
 \one_{(d,e)\le G}
 =\sum_{\substack{v\mid d\\v\mid e}}\lambda_G(v).}
 \label{eq:tpc33-row-gcd-projector}
\end{equation}
If \(e\) is squarefree, then
\begin{equation}
 \boxed{
 \sum_{v\mid e}|\lambda_G(v)|
 \le3^{\omega(e)}\ll_\eps e^\eps.}
 \label{eq:tpc33-row-gcd-projector-mass}
\end{equation}
\end{lemma}

\begin{proof}
Put \(n=(d,e)\).  Reversing the two divisor sums gives
\begin{align*}
 \sum_{v\mid n}\lambda_G(v)
 &=\sum_{\substack{g\mid n\\g\le G}}
   \sum_{w\mid n/g}\mu(w)\\
 &=\one_{n\le G}.
\end{align*}
This proves \eqref{eq:tpc33-row-gcd-projector}.  If \(e\) is squarefree,
then
\begin{align*}
 \sum_{v\mid e}|\lambda_G(v)|
 &\le\sum_{v\mid e}\sum_{g\mid v}|\mu(v/g)|
 =\sum_{gw\mid e}1=3^{\omega(e)}.
\end{align*}
The final bound is the standard divisor estimate.
\end{proof}

The key point is the order of use.  After the opposite row is fixed,
\(v\) ranges only over divisors of its squarefree coordinate \(e\), so the
total coefficient mass is \(X^{o(1)}\).  Summing
\eqref{eq:tpc33-row-gcd-projector} without fixing a column would not by
itself prove a small operator norm for the full mask.

\subsection{Reflecting one ultra leg}

Fix a right row
\begin{equation}
 \gamma=(\ell',e),\qquad n=m_\gamma=\ell'e,
 \label{eq:tpc33-fixed-right-column}
\end{equation}
and consider the left-ultra column
\begin{equation}
 B_\gamma^L
 :=\sum_\alpha\gamma_\alpha^{(1)}
 \sum_j\mathfrak A_{\alpha,\gamma}^{\rm act}(j)
 C_{m_\alpha}(j)\mathsf H_{m_\gamma}(j).
 \label{eq:tpc33-left-column-predefinition}
\end{equation}
Open the ultra increment.  For a nonzero term with
\(m_\alpha=\ell d\), let \(u\) be its ultra divisor and put
\begin{equation}
 N_\alpha(j)=\ell dj+h=us.
 \label{eq:tpc33-ultra-reflection}
\end{equation}
Then
\begin{equation}
 T<u\le U_0,\qquad
 1\le s<\frac{N_\alpha(j)}T\ll_{\mathscr D}\frac XT.
 \label{eq:tpc33-reflected-range}
\end{equation}
The variable \(s\) is not required to be squarefree.  Target primitivity
\eqref{eq:tpc33-divisor-primitivity} gives
\begin{equation}
 (s,\ell djH)=(u,\ell djH)=1.
 \label{eq:tpc33-reflected-primitivity}
\end{equation}
It does not imply \((s,u)=1\).

Fix \(\ell,j,s\).  Since \((s,\ell j)=1\), there is one residue
\(d_0\pmod s\) satisfying
\begin{equation}
 \ell jd_0+h\equiv0\pmod s.
 \label{eq:tpc33-d-residue}
\end{equation}
Writing \(d=d_0+s\tau\) gives
\begin{equation}
 u=u_0+\ell j\tau,\qquad
 u_0=\frac{\ell jd_0+h}{s},\qquad
 su_0-\ell jd_0=h.
 \label{eq:tpc33-unprojected-affine-column}
\end{equation}
The physical dyadic restriction \(d\asymp D\) permits at most
\begin{equation}
 O\!\left(1+\frac Ds\right)
 \label{eq:tpc33-unprojected-fiber-length}
\end{equation}
values of \(\tau\).

\subsection{The full actual-mask column theorem}

The row-gcd projector introduces the divisor \(v\).  On a nonempty term,
\(v\mid d,e\), and \eqref{eq:tpc33-reflected-primitivity} implies
\((v,s)=1\).  Write
\begin{equation}
 d=v\widetilde d.
 \label{eq:tpc33-v-projected-row}
\end{equation}
Since \((s,\ell jv)=1\), there is a unique
\(\widetilde d_0\pmod s\) such that
\begin{equation}
 \ell jv\widetilde d_0+h\equiv0\pmod s.
 \label{eq:tpc33-projected-residue}
\end{equation}
Put
\begin{equation}
 \widetilde u_0=\frac{\ell jv\widetilde d_0+h}{s}.
 \label{eq:tpc33-projected-u0}
\end{equation}

\begin{theorem}[Actual-mask affine-column disintegration]
\label{thm:tpc33-actual-mask-column}
Fix the right row \(\gamma=(\ell',e)\) in
\eqref{eq:tpc33-left-column-predefinition}.  After opening the left ultra
leg, the actual mask decomposes that column into the variables
\(\ell,j,s,v\), with \(v\mid e\), and at most two intervals \(I\) on each
such layer.  The inner arithmetic sequence has the exact form
\begin{equation}
 \boxed{
 \sum_{r\in I}
 \mu(\widetilde d_0+sr)
 \mu(\widetilde u_0+\ell jv r)
 \one_{(v,\, (\widetilde d_0+sr)
                 (\widetilde u_0+\ell jv r))=1}
 \rho(r)W(r),}
 \label{eq:tpc33-standard-affine-column}
\end{equation}
where
\begin{align}
 s\widetilde u_0-\ell jv\widetilde d_0&=h,
 \label{eq:tpc33-standard-affine-determinant}\\
 |I|&\ll_{\mathscr D}1+\frac{D}{sv}.
 \label{eq:tpc33-standard-affine-length}
\end{align}
Here \(\rho\) is a bounded fixed-period factor.  The weight \(W(r)\)
includes the varying factor \(\log u_r\), the row smooth factors, and their
Mellin/BV separations inherited from TPC-25; its seminorm cost is therefore
at most logarithmic in addition to the inherited soft losses.  The
coefficient outside \eqref{eq:tpc33-standard-affine-column} contains
\(-\mu(v)\lambda_G(v)(\log\ell)\) and the remaining fixed physical
factors.  Uniformly in the fixed column,
\begin{equation}
 \sum_{v\mid e}|\lambda_G(v)|\ll_\eps X^\eps.
 \label{eq:tpc33-column-mask-complexity}
\end{equation}
The local coprimality indicator in
\eqref{eq:tpc33-standard-affine-column} itself has a divisor expansion of
total coefficient mass at most \(2^{\omega(v)}\ll_\eps X^\eps\).
\end{theorem}

\begin{proof}
Insert \eqref{eq:tpc33-row-gcd-projector} into the last factor of
\eqref{eq:tpc33-actual-mask}.  Because \(e\) is fixed and squarefree,
\eqref{eq:tpc33-row-gcd-projector-mass} proves
\eqref{eq:tpc33-column-mask-complexity}.  For one divisor \(v\mid e\),
write \(d=v\widetilde d\) and use
\eqref{eq:tpc33-projected-residue}--
\eqref{eq:tpc33-projected-u0}.  Every solution is uniquely
\begin{equation}
 \widetilde d_r=\widetilde d_0+sr,
 \qquad
 u_r=\widetilde u_0+\ell jv r,
 \label{eq:tpc33-projected-affine-forms}
\end{equation}
and these forms satisfy
\eqref{eq:tpc33-standard-affine-determinant}.  The interval
\(d=v\widetilde d\asymp D\) has length \(O(D/v)\); sampling it with step
\(s\) proves \eqref{eq:tpc33-standard-affine-length}.

For fixed \(\ell,\ell',e,v\), the same-source condition is constant.  The
near-diagonal deletion is
\[
 |\ell v\widetilde d-\ell'e|>QX^{-\kappa_0}.
\]
Its complement is one interval on the \(\widetilde d\)-line, so the
retained set is the union of at most two intervals.  The fixed residue
factor is periodic with a modulus depending only on \(h\), and the smooth
 factors have the Mellin/BV separations already recorded in TPC-25.  These
 give \(\rho\) and \(W\) without an additional polynomial projective cost,
 apart from the logarithmic and soft losses recorded above.

The original row is squarefree, so \(v\) and \(\widetilde d_r\) are
squarefree and coprime.  The ultra coefficient is supported on squarefree
\(u_r\), and target primitivity makes \((v,u_r)=1\).  Therefore
\begin{equation}
 \mu(d)\bigl(-\mu(u_r)\log u_r\bigr)
 =-\mu(v)\mu(\widetilde d_r)\mu(u_r)\log u_r
 \label{eq:tpc33-projected-sign-fusion}
\end{equation}
on the physical support.  Restoring the coprimality by the displayed
indicator yields \eqref{eq:tpc33-standard-affine-column}.  Finally,
\[
 \one_{(v,n)=1}=\sum_{c\mid(v,n)}\mu(c)
\]
has coefficient mass \(2^{\omega(v)}\), which is \(X^{o(1)}\) because
\(v\mid e\ll X^{O(1)}\).
\end{proof}

\begin{corollary}[Both polarizations]
\label{cor:tpc33-both-column-disintegrations}
The right-ultra column obtained by fixing \(\alpha\) and opening
\(C_{m_\gamma}\) has the symmetric disintegration
\eqref{eq:tpc33-standard-affine-column}, with the two row labels exchanged.
Thus both terms in \eqref{eq:tpc33-matched-polarization} are covered.
\end{corollary}

\begin{remark}[Column complexity versus matrix complexity]
The theorem fixes the opposite row before expanding the gcd mask.  It does
not give an \(X^{o(1)}\) projective decomposition of the full matrix
\((\mathfrak m_{\rm act}(\alpha,\gamma))\).  In particular, it cannot be
substituted into an abstract Schur-multiplier argument without retaining
the outer column energy.
\end{remark}

\begin{remark}[Fixed determinant, growing slopes]
The two slopes in \eqref{eq:tpc33-standard-affine-column} are
\(s\) and \(a=\ell jv\).  Both are coprime to \(h\) and to one another,
and their determinant is the fixed integer \(h\).  They nevertheless grow
with \(X\), and the interval may contain only \(O(1)\) points when
\(sv>D\).  A theorem for one fixed affine pair cannot simply be summed over
these layers.
\end{remark}
